\documentclass{report}
\usepackage{amsmath,amssymb}
\setlength{\parindent}{0mm}
\setlength{\parskip}{1em}
\begin{document}
\begin{center}
\rule{6in}{1pt} \
{\large Rob Haelterman \\
{\bf Does Anderson always Accelerate Picard ?}}

Royal Military Academy (MWMW) \\ Renaissancelaan 30 \\ B-1000 Brussels \\ Belgium
\\
{\tt robby.haelterman@mil.be}\\
Dirk Van Eester\\
Joris Degroote\\
Silviu Cracana\end{center}

Partitioned solution problems have often been solved using the Picard
fixed-point iteration because of its simplicity. More recently, Anderson
Acceleration has been re-discovered as a sure means to accelerate the
Picard iteration. Indeed, it has been proven that, whenever the Picard
iteration converges, the Anderson acceleration technique will result in
equal or faster convergence. However, as we will show, Anderson
acceleration has not always been applied wisely to the Picard iteration,
as the fine print of the aforementioned proof is more often than not
overlooked.\\
We will show how Anderson can be applied to a convergent Picard
iteration, resulting in far worse convergence and even in divergence. We
will also explain the reasons for this surprising behavior and present a
sure way to avoid it.\\
For this purpose we use a simple fabricated pathological example together
with a more elaborate model describing the elementary physics of a plasma
inside a Tokamak.


\end{document}
